\subsection{函数演算中的代入}

采用上面引入的记号，I, p.~66 的推论 1 与 I, p.~67 的推论 2 的结论分别成为

\[
\chi (g (\boldsymbol {a})) = g \left(\chi \left(a _ {1}\right), \dots , \chi \left(a _ {n}\right)\right), \quad \operatorname{Sp} (g (\boldsymbol {a})) = g \left(\operatorname{Sp} ^ {n} (\boldsymbol {a})\right)
\]

其中 \(f \in \mathcal{O}(\mathrm{Sp}^n(\boldsymbol{a});\mathrm{A})\)，\(\chi \in \mathrm{X}(\mathrm{A})\)，\(g \in \mathcal{O}(\mathrm{Sp}^n(\boldsymbol{a}))\)。

现在我们来证明一个更一般的代入性质。

定理 4. --- 设 A 是一个交换幺复 Banach 代数，设 \(n \geqslant 1\) 是一个整数，且 \(\boldsymbol{a} = (a_1, \ldots, a_n) \in \mathrm{A}^n\)。设 \(\boldsymbol{f} = (f_1, \ldots, f_p)\)，其中 \(f_1, \ldots, f_p\) 是 \(\mathcal{O}(\mathrm{Sp}^n(\boldsymbol{a}))\) 的元素。\(\mathrm{Sp}^n(\boldsymbol{a})\) 在映射 \(\boldsymbol{z} \mapsto \boldsymbol{f}(\boldsymbol{z}) = (f_1(\boldsymbol{z}), \ldots, f_p(\boldsymbol{z}))\) 下的像等于 \(\mathrm{Sp}^p(\boldsymbol{f}(\boldsymbol{a}))\)。

对每个 \(g \in \mathcal{O}(\mathrm{Sp}^{p}(\boldsymbol{f}(\boldsymbol{a}));\mathrm{A})\)，复合芽 \(g \circ f\) 是 \(\mathcal{O}(\mathrm{Sp}^{n}(\boldsymbol{a});\mathrm{A})\) 的元素，并且有 \(g(\boldsymbol{f}(\boldsymbol{a})) = (g \circ \boldsymbol{f})(\boldsymbol{a})\)。

关于 \(\mathrm{Sp}^{n}(\boldsymbol{a})\) 的像的第一个断言由 I, p.~67 的推论 2 得出。为证明第二个断言，我们将使用下述引理。

引理 12. --- 设 K 是 \(\mathrm{Sp}^p (\pmb {f}(\pmb {a}))\) 的多项式凸包。我们有 \(g(\pmb {f}(\pmb {a})) = (g\circ \pmb {f})(\pmb {a})\)（对每个芽 \(g\in \mathcal{O}(\mathrm{K};\mathrm{A})\)）。

设 \(\Psi\) 是从 \(\mathcal{O}(\mathrm{K};\mathrm{A})\) 到 A 的映射，满足 \(\Psi(g) = (g\circ f)(\boldsymbol{a})\)。它是一个连续的保单位元态射，满足 \(\Psi(z_j) = f_j(\boldsymbol{a})\)，其中 \(z_{j}\) 是第 \(j\) 个坐标函数（在 \(\mathbf{C}^p\) 上）的芽。因此，当 \(g\) 是多项式函数的芽时，有 \(\Psi(g) = g(\boldsymbol{f}(\boldsymbol{a}))\)。由 I, p.~68 的定理 2，此公式对所有 \(g\in \mathcal{O}(\mathrm{K};\mathrm{A})\) 仍然成立。

现在来证明该定理。设 V 是 \(\mathrm{Sp}^{p}(\boldsymbol{f}(\boldsymbol{a}))\) 的开邻域，\(\widetilde{g}\in\mathcal{O}(V;A)\) 是一个全纯函数，它在 \(\mathrm{Sp}^{p}(\boldsymbol{f}(\boldsymbol{a}))\) 的某邻域中的芽等于 g。设 \(b\in C^{p+q}\) 是 \((\boldsymbol{f}(\boldsymbol{a}),\mathrm{V})\) 的一个包络（I, p.~70 的引理 11），\(\pi\colon V\times C^{q}\to V\) 是典范投影。

设 \(\widetilde{f}_1, \ldots, \widetilde{f}_p\) 是以 \(f_1, \ldots, f_p\) 为芽的全纯函数，并设 U 是 \(\mathrm{Sp}^n(\boldsymbol{a})\) 的开邻域，使得 \((\widetilde{f}_1, \ldots, \widetilde{f}_p)(\mathrm{U}) \subset \mathrm{V}\)。设 \(\pi'\) 是从 \(\mathrm{U} \times \mathbf{C}^q\) 到 U 上的典范投影。用 \(z_{n+1}, \ldots, z_{n+q}\) 表示最后 \(q\) 个坐标函数（在 \(\mathbf{C}^{n+q}\) 上）。记 \(h = \widetilde{g} \circ (\widetilde{f}_1, \ldots, \widetilde{f}_p)\)，并令

\[
\boldsymbol {c} = (a _ {1}, \dots , a _ {n}, b _ {p + 1}, \dots , b _ {p + q}) \in \mathrm{A} ^ {n + q}.
\]

映射 \(\widetilde{g} \circ \pi\) 在开邻域 \(V \times C^q\)（即 \(Sp^{p+q}(b)\) 的多项式凸包 L 的开邻域）中全纯。将引理 12 应用于 \(c\) 以及下列函数在 L 的邻域中的芽

\[
(\widetilde {f} _ {1} \circ \pi^ {\prime}, \ldots , \widetilde {f} _ {p} \circ \pi^ {\prime}, z _ {n + 1}, \ldots , z _ {n + q}),
\]

以及 \(\widetilde{g}\circ\pi\) 的芽，便得

\[
(g \circ \pi) \left((f _ {1} \circ \pi^ {\prime}) (\boldsymbol {c}), \dots , (f _ {p} \circ \pi^ {\prime}) (\boldsymbol {c}), z _ {n + 1} (\boldsymbol {c}), \dots , z _ {n + q} (\boldsymbol {c})\right) = (h \circ \pi^ {\prime}) (\boldsymbol {c}).
\]

由于 \(\pi'(\boldsymbol{c}) = \boldsymbol{a}\)，有 \((h \circ \pi')(\boldsymbol{c}) = h(\boldsymbol{a})\) 与 \((f_i \circ \pi')(\boldsymbol{c}) = f_i(\boldsymbol{a})\)（对 \(1 \leqslant i \leqslant p\)；全纯函数演算的性质 (CF3)）。此外，有 \(z_{n+j}(\boldsymbol{c}) = b_{p+j}\)（对 \(1 \leqslant j \leqslant q\)）。由性质 (CF2)，有

\[
(g \circ \pi) (f _ {1} (\boldsymbol {a}), \dots , f _ {p} (\boldsymbol {a}), b _ {p + 1}, \dots , b _ {p + q}) = h (\boldsymbol {a}),
\]

再应用一次性质 (CF3)，便推得 \(g(f_{1}(\boldsymbol{a}), \ldots, f_{p}(\boldsymbol{a})) = h(\boldsymbol{a})\)。

